\begin{explanationbox}

	\textbf{\textcolor{CExplanation}{一句话直觉}}

	数列的通项 $a_n$ 是前 $n$ 项和 $S_n$ 的逐项差，就像从累积中取增量。

	\medskip
	\textbf{\textcolor{CExplanation}{核心拆解}}
	\begin{description}[style=nextline, leftmargin=2.8em, labelsep=0.5em, itemsep=0.45em]
		\item[\textbf{要点一（分段求通项）：}] 当 $n=1$ 时直接取 $a_1=S_1$；当 $n\ge2$ 时，利用 $a_n=S_n-S_{n-1}$ 计算。
		\item[\textbf{要点二（验证合并条件）：}] 若由 $n\ge2$ 公式代入 $n=1$ 所得值恰好等于 $S_1$，则通项可写为统一表达式；否则必须写成分段形式。
	\end{description}

	\medskip
	\textbf{\textcolor{CExplanation}{几何本质（累积与增量）}}

	无直观几何背景，可理解为：前 $n$ 项和是前 $n-1$ 项和加上第 $n$ 项，故而第 $n$ 项是相邻和的差。

	\medskip
	\textbf{\textcolor{CExplanation}{代数意义（差分定义）}}

	$a_n=S_n-S_{n-1}$ 是离散的差分形式，与连续函数求导对应（求导是差分的极限），体现离散求和与差分的互逆关系。

	\medskip
	\textbf{\textcolor{CExplanation}{考点价值（求通项基础）}}
	\begin{itemize}[leftmargin=*, itemsep=0.5em, topsep=0.4em]
		\item \textbf{考法一（直接求通项）：} 已知 $S_n$ 表达式，直接套用公式求 $a_n$。
		\item \textbf{考法二（分段判断）：} 要求验证 $n=1$ 是否满足，从而决定是否分段。
	\end{itemize}

	\medskip
	\textbf{\textcolor{CExplanation}{顿悟点}}

	\textbf{抓住“差”字：通项是前 $n$ 项和的增量，但起始值需单独校核。}

	\medskip
	\textbf{\textcolor{CExplanation}{使用场景}}
	\begin{itemize}[leftmargin=*, itemsep=0.5em, topsep=0.4em]
		\item \textbf{场景一（已知 $S_n$ 求 $a_n$）：} 适用范围：任意数列，在已知 $S_n$ 表达式时可直接求通项。
		\item \textbf{场景二（验证 $n=1$ 一致性）：} 在得到 $n\ge2$ 表达式后，必须回代 $n=1$ 检查是否等于 $S_1$，以确定是否需要分段。
	\end{itemize}
\end{explanationbox}
